\chapter{Conclusion and Future Work}
This chapter concludes Phase~I and sets out the work planned for Phase~II. Section~7.1 summarises what was built and found. Section~7.2 describes the Phase-II extensions.

\section{Conclusion}
This report has presented Phase~I of an adaptive study-planning system for self-directed learners. The system estimates a learner's pace with hierarchical Bayesian calibration and a context-aware enriched-shrinkage model, detects behavioural shifts with CUSUM, projects the completion date with a Gaussian process and a cold-start composite, and regenerates the plan with a capacity-based booking model. These components are realised in a working, local-first web application, with the pace calibration and change detection served by a Python intelligence service and the projection and booking engine running in the browser.

Each component was compared offline against established baselines on synthetic data with known ground truth, using held-out learner archetypes, bootstrap confidence intervals, and multiple-comparison correction. The evaluation was reported honestly rather than selectively. Pace calibration is an honest null on retrospective recovery error, where the structured Bayesian variants do not separate from the incumbent, but the enriched-shrinkage calibrator is a Holm-surviving improvement on the prospective, short-history prediction the planner depends on. Change detection is a robust null: no candidate beats the CUSUM frontier, and the fastest detector, CUSUM, is the one shipped because detection feeds a human-confirmed prompt. Progress projection produced the clearest finding, that the raw Gaussian-process interval under-covers severely and that a split-conformal correction restores near-nominal coverage. Schedule generation is a genuine win for a dynamic-programming allocator that the system deliberately does not ship, because the day-by-day packing problem it solved was removed by the move to a booking model.

The contribution of Phase~I is an integrated, open-loop adaptive planner that extends the predict-then-optimize architecture~\cite{islam2024} into a loop over a single learner's own logged behaviour, together with a rigorous and honest account of which components improve on their baselines, which do not, and which improvements are validated but not yet shipped.

\section{Future Work}
Phase~II will close the loop and add assessment-based verification of genuine learning, which is the principal novelty of the project. Three lines of work follow.

The first is verified learning. The pace signal that drives the plan is currently taken from self-reported study time, which is known to be gameable~\cite{gaming2026}. Phase~II will generate concept-tagged assessments from the learner's own study material using large language models~\cite{itemgen2024,lecturequiz2026}, tag each item to a knowledge component~\cite{kctag2024}, grade responses automatically~\cite{llmgrade2024}, and estimate genuine concept mastery with cold-start knowledge tracing suited to a single learner with little history~\cite{clst2024,maml2026}. The verified-mastery estimate will be fed back into pace calibration, closing the loop shown dashed in Figure~\ref{fig:arch-revised}, so that the plan responds to evidence of learning rather than to logged time alone. Review timing will draw on spaced-repetition and personalised forgetting models~\cite{spaced2025,conceptkt2026}.

The second is to ship the refinements that Chapter~6 validated but that Phase~I did not deliver, in particular the split-conformal correction of the projection interval, which fixes the under-coverage of the raw Gaussian-process interval past the cold-start threshold.

The third is longitudinal validation on real data. Phase~I rests on synthetic ground truth; Phase~II will validate the system on a single learner's own logged sessions over time, comparing closed-loop adaptive planning against the static open-loop baseline on schedule adherence and measured learning gain.
